
\vspace{0.5em}
\textbf{Haus 3 – Partyhaus}\\
Im Gesamtszenario des Smart Beach Resort ergänzt Haus 3 die Funktionen von Haus 2, das den Gästezugang und den Check-in-Prozess automatisiert, um einen Use Case, der den Aufenthalt der Gäste durch gezielte Licht- und Audioinszenierungen erweitert. Haus 3 dient als Demonstrationsumgebung für eine ereignisgesteuerte Entertainment-Steuerung innerhalb des Resorts und zeigt, wie IoT-Knoten neben sicherheits- und verwaltungsorientierten Aufgaben auch erlebnisorientierte Dienste bereitstellen können.

Der Einzel-Use-Case basiert auf der Verwendung des im KS5009-Kit enthaltenen RGB-LED-Moduls sowie des passiven Buzzers, die gemeinsam eine abgestimmte Sequenz aus Licht- und Toneffekten erzeugen. Die Aktivierung dieser Szene erfolgt über das zentrale Node-RED-Dashboard, das dem Betreiber die Buchung einer „Party“ ermöglicht. Während eine solche Party aktiv ist, können Gäste außerdem mittels der im Projekt vorhandenen RFID-Chips unterschiedliche Musikstücke abrufen, wodurch das System eine Art „Tonibox für junggebliebene Urlaubende“ abbildet und interaktiv einige Musikstücke in die Entertainment-Logik integriert. Das LCD-Display des Hauses zeigt parallel den aktuellen Szenenstatus wie „Party laeuft“ oder den Zeitpunkt der nächsten, geplanten Feier und dient so als lokale Rückmeldung über den Betriebsmodus.

Aus architektonischer Sicht demonstriert Haus 3 damit die Service-Schicht innerhalb des Gesamtsystems. Während andere Häuser sicherheits- oder umweltbezogene Aufgaben übernehmen, repräsentiert dieses Haus den Aspekt der User Experience  im IoT-Kontext. Die Implementierung zeigt, dass verteilte Mikrocontroller nicht nur Messdaten erfassen, sondern auch ereignisgesteuert auf Benutzeranforderungen reagieren können.

Da bei einer Party laute Musik eine gängige Begleiterscheinung darstellt und die Häuser ein gemeinsames Resort bilden, wird beim Start des Partymodus ein Schutzmechanismus aktiviert: Die Fenster und Türen der übrigen Häuser werden automatisch geschlossen, um mögliche Lärmbelästigungen für andere Gäste zu minimieren.

